\begin{abstract}
Benchmarks for single-cell clustering usually fix each method's setting with the reference
labels, so their scores assume information a user does not have. We ask what happens to the
advantage scMINER reports for its mutual information clustering, MICA, over a principal
component and Leiden pipeline when one label-free rule chooses the setting for both. The rule
is subsampling stability in the style of chooseR, fixed before scoring and run at three seeds
on all ten benchmark datasets. On the nine small datasets MICA leads narrowly under three of
four rules, but its mean adjusted Rand index moves by 0.255 across rules against 0.120 for the
baseline, and two readings of the same stability scores reverse the ranking. The reversal comes
from the candidates. Capping both methods at 20 clusters does not give them the same cluster
counts, and MICA's two readings land at extremes the baseline cannot reach. Restricting MICA's
saved scores after the fact to the counts the baseline produced removes it, with MICA ahead
under both readings, though its finest-stable picks then often sit at the edge of that range, so
this shows what the reversal depends on rather than giving a clean ranking. The tie-break still
moves MICA more than the baseline, 0.12 against 0.07 per dataset. The same confound
explains an apparent difference between MICA's two pipelines run on the same 2,500 cells. On
PBMC14K, the one large benchmark dataset, both methods choose the coarsest partition their grid allows,
and a 0.02 change in the rule's threshold reverses which leads. We also compare a faster
neighbour graph, a correlation shortlist rescored exactly, with MICA's own search. On PBMC14K
it matches MICA's recall and clustering in a small fraction of the build time. On 100,000
Allen cells it takes less than half the build time and reaches higher recall, but gives no better
clustering. On both it needs three to four times the memory.
\end{abstract}

\begin{IEEEkeywords}
Single-cell RNA sequencing, mutual information clustering, label-free model selection,
cluster number selection, benchmark design
\end{IEEEkeywords}
